\section{\Glspl{vm}}

\Glspl{vm} são a correspondência em software --- em nível de virtualização --- de máquinas físicas, podendo conter CPUs virtuais, memória e dispositivos periféricos.
Em correspondência com \glspl{so} modernos, \glspl{vm} funcionam de forma análoga a processos no projeto. Por definição, a criação de uma \gls{vm} deve receber uma porção
de memória alocada para si; esta será a região delimitada e autorizada a ser acessada pelo convidado. Além disso, para cada núcleo esperado pelo convidado, a \gls{vm} criará uma CPU virtual,
ficando também responsável por criar e gerenciar os dispositivos virtuais.

\begin{figure}[H]
\centering
\includegraphics[width=0.8\textwidth]{images/SimplifiedVirtualMachine.pdf}
\caption{Diagrama de Classes Simplificado de \glspl{vm}}\label{fig:simplified-virtual-machine}
\end{figure}
